\documentclass[a4paper,12pt]{article}
\usepackage[utf8]{inputenc}
\usepackage[T2A]{fontenc}
\usepackage[russian,english]{babel}
\usepackage{geometry}
\geometry{left=2.5cm,right=1.5cm,top=2cm,bottom=2cm}
\usepackage{graphicx}
\usepackage{hyperref}
\usepackage{verbatim}
\usepackage{indentfirst}

\begin{document}

% --- TITLE PAGE ---


\begin{center}

{\large
Министерство науки и высшего образования\\
Российской Федерации
}

\vspace{0.7cm}

{\large\bfseries
ФЕДЕРАЛЬНОЕ ГОСУДАРСТВЕННОЕ\\
АВТОНОМНОЕ ОБРАЗОВАТЕЛЬНОЕ\\
УЧРЕЖДЕНИЕ\\
ВЫСШЕГО ОБРАЗОВАНИЯ
}

\vspace{0.5cm}

{\large\bfseries
«НАЦИОНАЛЬНЫЙ ИССЛЕДОВАТЕЛЬСКИЙ\\
УНИВЕРСИТЕТ ИТМО»
}

\vspace{0.45cm}

{\normalsize
Факультет систем управления и робототехники
}

\vspace{2.8cm}

{\Large\bfseries
ЛАБОРАТОРНАЯ РАБОТА №0
}

\vspace{0.3cm}

{\normalsize
по дисциплине «Машинное обучение»
}

\vspace{0.8cm}

{\LARGE
Знакомство с git и GitHub
}

\vfill

\begin{flushright}
\begin{minipage}{0.43\textwidth}

\textbf{Выполнили:}

Студент: Нколе Эндрю,\\
\textbf{Группа: R3343}

\vspace{0.3cm}

\textbf{Преподаватель:}

 Поляк Марк Дмитриевич

\end{minipage}
\end{flushright}

\vfill

Санкт-Петербург\\
2026

\end{center}

\end{titlepage}


\tableofcontents
\newpage

% --- TASK TEXT ---
\section{Текст задания}
\textbf{Цель:} поразмышлять о своих целях и задачах на текущий семестр. Создать файлы в репозитории, освоить работу с Git из командной строки, подготовить отчет.

\textbf{Задание 1:} Создать файл goals.md, добавить заголовки, текст о своем опыте, нумерованный и ненумерованный списки, создать папку resources с изображением.

\textbf{Задание 2:} Клонировать репозиторий, создать info.md, выполнить коммиты, создать .gitignore и объяснить результат игнорирования файла data/dataset.csv.

\textbf{Задание 3:} Подготовить отчет в формате PDF и загрузить его в репозиторий.

% --- TASK 1 ---
\section{Результат выполнения Задания 1}
\subsection{Содержимое файла goals.md}
\begin{verbatim}
# Мой опыт

Я начал изучать программирование в университете. До этого я немного программировал на Python и делал простые студенческие проекты. Мой профиль на GitHub: Andy13080(https://github.com/Andy13080/git-lab.git). Теперь я хочу научиться лучше работать в команде и освоить Git.

# Мои цели

1. Научиться уверенно пользоваться Git и GitHub.

2. Научиться готовить отчёты в LaTeX.

3. Научиться управлять проектом в репозитории.

## Задачи

- Изучить основные команды Git: add, commit, push, pull.

- Понять, как работает .gitignore.

- Подготовить отчёт в LaTeX и загрузить его в репозиторий.

- Научиться форматировать документы в Markdown.

- Регулярно делать коммиты с понятными сообщениями.
\end{verbatim}
\textit{Примечание: Изображение было успешно загружено в папку resources и добавлено в файл goals.md.}

% --- TASK 2 ---
\section{Результат выполнения Задания 2}
\subsection{Содержимое файла info.md}
\begin{verbatim}
НКОЛЕ ЭНДРЮ

Windows 11 Home

2026-09-18 21:20
\end{verbatim}

\subsection{Содержимое файла .gitignore}
\begin{verbatim}
data/
\end{verbatim}

\subsection{Объяснение результата попытки закоммитить файл data/dataset.csv}
В файл .gitignore была добавлена строка \texttt{data/}. Это указывает Git игнорировать всё содержимое папки data. Поэтому при выполнении команды \texttt{git add data/dataset.csv} Git проигнорировал этот файл, и он не был добавлен в индекс. Соответственно, коммит с этим файлом не создался. Это доказывает, что .gitignore работает корректно.

% --- CONCLUSION ---
\section{Вывод}
В ходе выполнения данной лабораторной работы я ознакомился с системой контроля версий Git и сервисом GitHub. Я научился создавать репозитории, работать с файлами через веб-интерфейс и командную строку, а также освоил базовые команды Git (add, commit, push) и настройку .gitignore.

\end{document}